% Generated by scripts/assemble.py; corrections belong in source/corrections.json.
\examhead{30}{中国石油大学(北京) 2025 年数学分析试卷}{33}
% p034-b19; source PDF page 34
\par\Needspace{5\baselineskip}
\blockanchor{p034-b19}
\majorquestion{一}计算题. 每题 10 分, 共 60 分.\par
% p034-b20; source PDF page 34
\blockanchor{p034-b20}
\question{1}设 $x_{n} = \cos \frac{x}{2} \cos \frac{x}{2^{2}} \cdots \cos \frac{x}{2^{n}}$，求 $\lim_{x \to 0} \lim_{n \to \infty} x_{n}$.\par
% p034-b21; source PDF page 34
\blockanchor{p034-b21}
\question{2}求极限 $\lim_{n\to \infty}\frac{\sqrt[n]{n!}}{n}$.\par
% p035-b00; source PDF page 35
\blockanchor{p035-b00}
\question{3}已知 $f(x) = \begin{cases} x^2 \sin \frac{\pi}{x}, & x > 0; \\ A, & x = 0; \\ ax^2 + b, & x < 0. \end{cases}$$f(x)$ 在 $x = 0$ 处可导，求 $A, a, b$.\par
% p035-b01; source PDF page 35
\blockanchor{p035-b01}
\question{4}$a_{1}, a_{2}, \cdots, a_{n}$  为正数, n > 2,  $f(x) = \left(\frac{a_{1}^{x} + a_{2}^{x} + \cdots + a_{n}^{x}}{n}\right)^{\frac{1}{x}}$ , 求  $\lim_{x \to 0^{+}} f(x)$ .\par
% p035-b02; source PDF page 35
\blockanchor{p035-b02}
\question{5}求不定积分 $\int \frac{1}{x} \sqrt{\frac{x + 2}{x - 2}} \, \mathrm{d}x$.\par
% p035-b03; source PDF page 35
\blockanchor{p035-b03}
\question{6}求 $I = \int_{C} (x^{2} + 2xy) \, \mathrm{d}y$, $C$ 表示逆时针方向上半椭圆 $\frac{x^{2}}{a^{2}} + \frac{y^{2}}{b^{2}} = 1$.\par
% p035-b04; source PDF page 35
\par\Needspace{5\baselineskip}
\blockanchor{p035-b04}
\majorquestion{二}(15 分) 计算 $f(x) = \arcsin x$ 的麦克劳林公式到 $x^5$ 项.\par
% p035-b05; source PDF page 35
\par\Needspace{5\baselineskip}
\blockanchor{p035-b05}
\majorquestion{三}(15 分) 求 $f(x, y) = x^2 - xy + y^2 - 2x + y$ 在全平面上的最大与最小值.\par
% p035-b06; source PDF page 35
\par\Needspace{5\baselineskip}
\blockanchor{p035-b06}
\majorquestion{四}(15 分) 作极坐标变换, 将二重积分 $\iint_{D} f\left(\sqrt{x^2 + y^2}\right) \mathrm{d}x \mathrm{~d}y$ 为定积分, $D = \{(x, y) \mid 0 \leq y \leq x \leq 1\}$.\par
% p035-b07; source PDF page 35
\par\Needspace{5\baselineskip}
\par\Needspace{6\baselineskip}
\blockanchor{p035-b07}
\majorquestion{五}(15 分) 叙述数列 $\{x_{n}\}$ 收敛的柯西收敛原理, 并利用它证明: 若数列 $\{b_{n}\}$ 有界, 则\par
\nopagebreak[4]
% p035-b08; source PDF page 35
\blockanchor{p035-b08}
\displayformula{a _ {n} = \frac {b _ {1}}{1 \cdot 2} + \frac {b _ {2}}{2 \cdot 3} + \dots + \frac {b _ {n}}{n (n + 1)} (n \in \mathbb {N} ^ {+})}
% p035-b09; source PDF page 35
\blockanchor{p035-b09}
收敛.\par
% p035-b10; source PDF page 35
\par\Needspace{5\baselineskip}
\blockanchor{p035-b10}
\majorquestion{六}(15 分) 设 $a_{n} > 0$ 且单调递减, 证明: $\sum_{n=1}^{\infty} a_{n}$ 与 $\sum_{n=1}^{\infty} 2^{n} a_{2^{n}}$ 同敛散.\par
% p035-b11; source PDF page 35
\par\Needspace{5\baselineskip}
\blockanchor{p035-b11}
\majorquestion{七}(15 分) 讨论积分 $I(p) = \int_{0}^{1} x^{p-1} (\ln x)^{2} \, \mathrm{d}x$ 在 $[p_0, +\infty)$ 上的一致收敛性, $p_0 > 0$.\par
